\subsection{\textit{Instance Segmentation}}

\textit{Instance segmentation} merupakan tugas \textit{Computer Vision} yang menghasilkan informasi kelas dan wilayah piksel untuk setiap instans objek. Berbeda dengan segmentasi semantik yang memberikan label kelas pada piksel tanpa membedakan objek dari kelas yang sama, \textit{instance segmentation} mempertahankan identitas setiap objek sebagai instans yang terpisah \cite{he2017maskrcnn}.

Pemisahan pada tingkat instans sesuai dengan kebutuhan sistem karena satu piring dapat mengandung lebih dari satu lauk dari kelas yang sama. Setiap objek perlu dipertahankan sebagai instans yang berbeda agar jumlahnya dapat dihitung secara individual. Pemanfaatan segmentasi dan penghitungan pada tingkat instans juga digunakan pada penelitian SibNet dan FoodMask untuk mengenali serta menghitung objek makanan \cite{nguyen2022sibnet,nguyen2023foodmask}.

Kondisi oklusi dapat menjadi tantangan dalam \textit{instance segmentation} karena sebagian area objek dapat tertutup oleh objek lain sehingga informasi visual yang tersedia berkurang \cite{ke2021}. Oleh karena itu, kondisi oklusi dipandang sebagai salah satu faktor yang dapat memengaruhi hasil prediksi model dan bukan sebagai kondisi yang diasumsikan selalu dapat ditangani secara sempurna.

Pada sistem yang dirancang, keluaran \textit{instance segmentation} digunakan untuk memperoleh kelas dan identitas setiap instans lauk. Luas \textit{segmentation mask} tidak digunakan untuk menentukan harga maupun jumlah lauk; jumlah ditentukan berdasarkan banyaknya instans yang dipertahankan sebagai hasil prediksi akhir.
